\documentclass[12pt]{article}
\usepackage[utf8]{inputenc}
\usepackage[spanish]{babel}
\usepackage{amsmath}
\usepackage{amssymb}
\usepackage{booktabs}
\usepackage{geometry}
\geometry{margin=2.5cm}
\usepackage[most]{tcolorbox}
\newtcolorbox{intuicion}{
  colback=blue!5!white,
  colframe=blue!40!black,
  fonttitle=\bfseries,
  title=Intuici\'on,
  boxrule=0.6pt,
  arc=4pt,
  left=6pt, right=6pt, top=4pt, bottom=4pt
}
\newtcolorbox{intuicionmat}{
  colback=green!5!white,
  colframe=green!40!black,
  fonttitle=\bfseries,
  title=\textit{?`Qu\'e dice esta ecuaci\'on?},
  boxrule=0.6pt,
  arc=4pt,
  left=6pt, right=6pt, top=4pt, bottom=4pt
}
\newtcolorbox{explicacion}{
  colback=orange!5!white,
  colframe=orange!50!black,
  fonttitle=\bfseries,
  title=Explicaci\'on,
  boxrule=0.6pt,
  arc=4pt,
  left=6pt, right=6pt, top=4pt, bottom=4pt
}

\title{\textbf{Resumen: Cap\'itulo 5 --- Paseos Aleatorios y Movimiento Browniano}\\
\large \textit{(Random Walks and Brownian Motion)}\\
\normalsize J.~Robert~Buchanan, \textit{An Undergraduate Introduction to Financial Mathematics}, 2nd~ed.}
\author{}
\date{Junio 2026}

\begin{document}
\maketitle
\tableofcontents
\newpage

%%% =========================================================
\section{Idea Intuitiva del Paseo Aleatorio (Intuitive Idea of a Random Walk)}
%%% =========================================================

\subsection{Motivaci\'on y definici\'on informal}

El cap\'itulo modela el comportamiento estoc\'astico de precios de acciones, opciones
e \'indices mediante \textbf{paseos aleatorios (random walks)}.
La idea central proviene de un experimento simple: una persona parada en el origen
de la recta real lanza una moneda repetidamente. Si sale cara avanza una unidad
hacia la derecha; si sale cruz retrocede una unidad hacia la izquierda.
La trayectoria resultante es una realizaci\'on del paseo aleatorio.

\begin{intuicion}
Un paseo aleatorio es la versi\'on matem\'atica de un jugador que apuesta una moneda
una y otra vez. Su posici\'on va cambiando de forma impredecible,
pero existe una estructura probabil\'istica precisa que gobierna cu\'ales posiciones
puede alcanzar y con qu\'e probabilidad.
\end{intuicion}

Generalizaciones importantes:
\begin{itemize}
  \item La moneda no necesita ser justa ($p \neq 1/2$).
  \item El tama\~no del paso puede variar en cada etapa.
  \item El desplazamiento puede provenir de una distribuci\'on normal.
\end{itemize}

%%% =========================================================
\section{An\'alisis del Primer Paso (First Step Analysis)}
%%% =========================================================

\subsection{Representaci\'on como suma parcial}

Sea $X_n$ una variable aleatoria que toma el valor $+1$ con probabilidad
$p = 1/2$ y el valor $-1$ con probabilidad $1/2$.
Las $X_i$ se asumen \textbf{independientes e id\'enticamente distribuidas (i.i.d.)}.
El \textbf{$N$-\'esimo estado del paseo (N-th state of the walk)} es la suma parcial:

\begin{equation}
  S(N) = S(0) + X_1 + X_2 + \cdots + X_N, \qquad N > 0.
  \label{eq:rw}
\end{equation}

\begin{intuicionmat}
Esta ecuaci\'on dice que la posici\'on actual del caminante es la posici\'on
inicial m\'as la suma acumulada de todos los pasos dados. Cada $X_k$ es
un salto aleatorio de $\pm 1$, de modo que $S(N)$ oscila alrededor del
punto de partida. En finanzas, $S(N)$ puede interpretarse como el precio
de un activo en el d\'ia $N$ si el precio cambia en $\pm 1$ cada d\'ia.
\end{intuicionmat}

\begin{intuicion}
La ecuaci\'on~\eqref{eq:rw} es la piedra angular del cap\'itulo: conecta la posici\'on
en el tiempo $N$ con todos los pasos previos. Sin ella no se puede calcular ninguna
probabilidad ni tiempo de parada.
\end{intuicion}

La definici\'on inductiva equivalente es $S(N) = S(N-1) + X_N$ para $N > 0$.
Sin p\'erdida de generalidad se toma $S(0) = 0$ (translaci\'on de variable
$T(n) = S(n) - S(0)$).

Si en $n$ pasos se dan $k \geq 0$ pasos positivos y $n-k \geq 0$ negativos,
la probabilidad de alcanzar el estado $S(n) = 2k - n$ es:

\begin{equation}
  \mathrm{P}\!\left(T(n) = 2k - n\right) = \binom{n}{k}\!\left(\frac{1}{2}\right)^{\!n}.
  \label{eq:binom}
\end{equation}

\begin{intuicionmat}
Esta es la f\'ormula binomial cl\'asica con $p = 1/2$. Indica que para llegar
a la posici\'on $2k-n$ se necesitan exactamente $k$ pasos hacia la derecha entre
$n$ pasos totales, y hay $\binom{n}{k}$ secuencias que logran eso, cada una con
probabilidad $(1/2)^n$.
\end{intuicionmat}

\subsection{Lema 5.1: Estados alcanzables}

\begin{tcolorbox}[colback=gray!5, colframe=gray!50, title=\textbf{Lema~5.1}]
Para el paseo aleatorio~\eqref{eq:rw} con estado inicial $S(0)=0$:
\begin{enumerate}
  \item $\mathrm{P}(S(n)=m)=0$ si $|m|>n$,
  \item $\mathrm{P}(S(n)=m)=0$ si $n+m$ es impar,
  \item $\mathrm{P}(S(n)=m)=\displaystyle\binom{n}{(n+m)/2}\!\left(\frac{1}{2}\right)^{\!n}$
        en caso contrario.
\end{enumerate}
\end{tcolorbox}

\begin{intuicion}
En $n$ pasos la posici\'on s\'olo puede cambiar como m\'aximo en $n$ unidades,
por lo que $|S(n)| \leq n$. Adem\'as, cada paso cambia la paridad de la posici\'on,
as\'i que despu\'es de $n$ pasos s\'olo son alcanzables estados con la misma paridad
que $n$.
\end{intuicion}

\subsection{Condici\'on de frontera absorbente (Absorbing boundary condition)}

En el modelo financiero el precio de un activo sigue un paseo aleatorio.
Si el precio cae a cero, el activo no tiene valor y el caminante \emph{queda atrapado}.
Formalmente:

\begin{itemize}
  \item Se define $m_n = \min\{S(k) : 0 \leq k \leq n\}$, el m\'inimo hist\'orico.
  \item Si $S(0) > 0$ y la frontera inferior es $0$, en cuanto $m_i = 0$ se
        impone $S(k) = 0$ para todo $k \geq i$.
\end{itemize}

Las dos preguntas centrales son:
\begin{enumerate}
  \item ?`Cu\'al es la probabilidad de cruzar un umbral $A > 0$ sin tocar la frontera $0$?
  \item ?`Cu\'al es el n\'umero esperado de pasos antes de cruzar $A > 0$
        por primera vez?
\end{enumerate}

\begin{intuicion}
La frontera absorbente en $0$ modela la quiebra: una vez que el precio de la acci\'on
llega a cero, el inversor pierde todo y el proceso se detiene. El umbral $A$
representa un nivel de ganancia objetivo.
\end{intuicion}

\subsection{Lema 5.2: F\'ormula de reflexi\'on}

\begin{tcolorbox}[colback=gray!5, colframe=gray!50, title=\textbf{Lema~5.2}]
Bajo las condiciones de~\eqref{eq:rw} con $X_i = \pm 1$ equidistribuidas
y frontera absorbente en $0$, para $A, i > 0$:
\begin{equation}
  \mathrm{P}(S(n) = A \wedge m_n > 0 \mid S(0) = i)
  = \left[\binom{n}{(n+A-i)/2} - \binom{n}{(n-A-i)/2}\right]
    \!\left(\frac{1}{2}\right)^{\!n}.
  \label{eq:reflec}
\end{equation}
\end{tcolorbox}

\textbf{Demostraci\'on (esquema).}
Se aplica la \textbf{propiedad de homogeneidad espacial (spatial homogeneity)}:
trasladar el origen transforma $\mathrm{P}(S(n)=A \mid S(0)=i)$ en
$\mathrm{P}(S(n)=A-i \mid S(0)=0)$.
Por \textbf{simetr\'ia del paseo no restringido (symmetry of the unrestricted walk)},
$\mathrm{P}(S(n)=j \mid S(0)=0) = \mathrm{P}(S(n)=-j \mid S(0)=0)$.
Cualquier trayectoria que toca la frontera $0$ puede \emph{reflejarse} en el eje,
de modo que su peso probabil\'istico es igual al de la trayectoria reflejada
que llega a $-A$ sin restricci\'on:
\begin{equation*}
  \mathrm{P}(S(n) = A \wedge m_n \leq 0 \mid S(0)=i)
  = \binom{n}{(n-A-i)/2}\!\left(\frac{1}{2}\right)^{\!n}.
\end{equation*}
Restando del total (sin restricci\'on) se obtiene~\eqref{eq:reflec}. \hfill $\square$

\begin{intuicionmat}
La ecuaci\'on~\eqref{eq:reflec} cuantifica la probabilidad de \emph{llegar a $A$
sin haber tocado $0$}. El primer t\'ermino cuenta todas las trayectorias que llegan
a $A$; el segundo substrae las que pasaron por $0$ en alg\'un momento
(usando el truco de reflejo: esas trayectorias tienen la misma probabilidad que
las que llegan a $-A$ sin restricci\'on).
\end{intuicionmat}

\subsection{Notaci\'on compacta y tiempo de parada}

Se define la funci\'on auxiliar:
\begin{equation}
  f_{A,i}(n) = \left[\binom{n}{(n+A-i)/2} - \binom{n}{(n-A-i)/2}\right]
               \!\left(\frac{1}{2}\right)^{\!n}.
\end{equation}

El \textbf{tiempo de parada (stopping time)} $\Omega_A$ es la primera vez que
$S(n) = A$. Se puede demostrar que:
\begin{equation}
  \mathrm{P}(\Omega_A = n \mid S(0) = i) = \tfrac{1}{2}\,f_{1,(i-A)}(n-1).
\end{equation}

\subsection{Ecuaci\'on de Laplace discreta y probabilidades de absorci\'on}

Suponga que $0 < S(0) = i < A$ con fronteras absorbentes en $0$ y en $A$.
Sea $\mathcal{P}_A(i)$ la probabilidad de que el paseo alcance $A$ evitando $0$.
Descomponiendo en el primer paso:

\begin{equation}
  0 = \mathcal{P}_A(i-1) - 2\mathcal{P}_A(i) + \mathcal{P}_A(i+1),
  \label{eq:laplace}
\end{equation}

con condiciones de frontera $\mathcal{P}_A(0) = 0$ y $\mathcal{P}_A(A) = 1$.

\begin{intuicionmat}
La ecuaci\'on~\eqref{eq:laplace} es la aproximaci\'on en diferencias finitas de la
segunda derivada igualada a cero: $\mathcal{P}_A''(i) \approx 0$. Se llama
\textbf{ecuaci\'on de Laplace discreta (discrete Laplacian equation)}.
Matem\'aticamente dice que la probabilidad de llegar a $A$ en el estado $i$
es el promedio de las probabilidades en los estados vecinos $i-1$ e $i+1$,
lo cual es exactamente lo que uno esperar\'ia si cada paso es igualmente
probable hacia la izquierda o la derecha.
\end{intuicionmat}

En forma matricial el sistema~\eqref{eq:laplace} con sus condiciones de frontera es:
\begin{equation*}
  \begin{bmatrix}
    1 & 0 & 0 & \cdots & 0 & 0 \\
    1 & -2 & 1 & \cdots & 0 & 0 \\
    0 & 1 & -2 & \cdots & 0 & 0 \\
    \vdots & & & \ddots & & \vdots \\
    0 & 0 & 0 & \cdots & -2 & 1 \\
    0 & 0 & 0 & \cdots & 0 & 1
  \end{bmatrix}
  \begin{bmatrix} \mathcal{P}_A(0) \\ \mathcal{P}_A(1) \\ \mathcal{P}_A(2) \\ \vdots \\
  \mathcal{P}_A(A-1) \\ \mathcal{P}_A(A) \end{bmatrix}
  =
  \begin{bmatrix} 0 \\ 0 \\ 0 \\ \vdots \\ 0 \\ 1 \end{bmatrix}.
\end{equation*}

\subsection{Teorema 5.1: Probabilidades de absorci\'on}

\begin{tcolorbox}[colback=gray!5, colframe=gray!50, title=\textbf{Teorema~5.1}]
Bajo las condiciones del Lema~5.2 con fronteras absorbentes en $0$ y en $A$,
si $0 \leq S(0) = i \leq A$:
\begin{enumerate}
  \item La probabilidad de alcanzar $A$ sin tocar $0$ es $\mathcal{P}_A(i) = i/A$.
  \item La probabilidad de alcanzar $0$ sin tocar $A$ es $\mathcal{P}_0(i) = 1 - i/A$.
\end{enumerate}
\end{tcolorbox}

\begin{intuicion}
El resultado es sorprendentemente elegante: la probabilidad de ganar (llegar a $A$)
es simplemente la fracci\'on de la distancia total que ya se ha recorrido.
Si el precio de una acci\'on comienza en $\$5$ y uno quiere ganar si llega a $\$10$
(o perder si llega a $\$0$), la probabilidad de ganar es exactamente $5/10 = 50\%$.
\end{intuicion}

\subsection{Ecuaci\'on de Poisson y tiempo esperado de salida}

Sea $B = \{0, A\}$ el conjunto de las dos fronteras absorbentes.
Define:
\begin{itemize}
  \item $\omega_{p_{i \to B}}$: tiempo de salida del camino $p_{i \to B}$ que parte de $S(0) = i$.
  \item $\Omega_B(i)$: valor esperado del tiempo de salida para $S(0) = i$.
\end{itemize}

Descomponiendo en el primer paso se obtiene la recursi\'on:

\begin{equation}
  -2 = \Omega_B(i-1) - 2\Omega_B(i) + \Omega_B(i+1),
  \label{eq:poisson}
\end{equation}

con $\Omega_B(0) = \Omega_B(A) = 0$.

\begin{intuicionmat}
La ecuaci\'on~\eqref{eq:poisson} es la \textbf{ecuaci\'on de Poisson discreta
(discrete Poisson equation)}: es id\'entica a la de Laplace pero con lado derecho
$-2$ en lugar de $0$. Ese $-2$ contabiliza el costo de un paso adicional
que toma el proceso en cada estado interior.
\end{intuicionmat}

En forma matricial:
\begin{equation}
  \begin{bmatrix}
    1 & 0 & \cdots & 0 & 0 \\
    1 & -2 & \cdots & 0 & 0 \\
    \vdots & & \ddots & & \vdots \\
    0 & 0 & \cdots & -2 & 1 \\
    0 & 0 & \cdots & 0 & 1
  \end{bmatrix}
  \begin{bmatrix} \Omega_B(0) \\ \Omega_B(1) \\ \vdots \\ \Omega_B(A-1) \\ \Omega_B(A) \end{bmatrix}
  =
  \begin{bmatrix} 0 \\ -2 \\ \vdots \\ -2 \\ 0 \end{bmatrix}.
  \label{eq:poissonmat}
\end{equation}

\subsection{Eliminaci\'on gaussiana y sustituci\'on hacia atr\'as}

Las matrices de los sistemas~\eqref{eq:laplace} y~\eqref{eq:poisson}
son \textbf{matrices tridiagonales (tridiagonal matrices)}: todos sus elementos son
cero excepto los de la diagonal principal, la subdiagonal y la superdiagonal.
Se resuelven mediante:

\begin{enumerate}
  \item \textbf{Eliminaci\'on gaussiana (Gaussian elimination)}: se suman
        m\'ultiplos de filas para convertir la \textbf{matriz aumentada (augmented matrix)}
        en forma \textbf{triangular superior (upper triangular)}.
        Tras aplicar el procedimiento por inducci\'on (multiplicando la fila $i$
        por $1 - 1/i$ y sum\'andola a la fila $i+1$), la matriz aumentada del
        sistema de Poisson queda:
        \begin{equation}
          \left[\begin{array}{ccccccc|c}
            1 & 0 & 0 & \cdots & 0 & 0 & 0 & 0 \\
            0 & -2 & 1 & \cdots & 0 & 0 & 0 & -2 \\
            0 & 0 & -\tfrac{3}{2} & \cdots & 0 & 0 & 0 & -3 \\
            \vdots & & & \ddots & & & & \vdots \\
            0 & 0 & 0 & \cdots & -\tfrac{A}{A-1} & 1 & 0 & -A \\
            0 & 0 & 0 & \cdots & 0 & 0 & 1 & 0
          \end{array}\right].
          \label{eq:upper}
        \end{equation}
  \item \textbf{Sustituci\'on hacia atr\'as (backwards substitution)}: conocido $\Omega_B(j)$
        se recupera $\Omega_B(j-1)$ mediante:
        \begin{equation}
          \Omega_B(j-1) = j - 1 + \frac{j-1}{j}\,\Omega_B(j).
          \label{eq:backsub}
        \end{equation}
\end{enumerate}

\begin{intuicionmat}
La ecuaci\'on~\eqref{eq:backsub} es la clave de la sustituci\'on hacia atr\'as:
dado que se conoce el tiempo esperado desde $j$, se puede calcular el tiempo
esperado desde $j-1$ sin resolver el sistema completo de nuevo.
\end{intuicionmat}

\begin{intuicion}
La eliminaci\'on gaussiana transforma un sistema de ecuaciones complicado en otro
equivalente pero en forma escalonada, que se resuelve f\'acilmente de abajo hacia
arriba. Es el m\'etodo est\'andar de \'algebra lineal num\'erica.
\end{intuicion}

\subsection{Teorema 5.2: Tiempo esperado de salida incondicional}

\begin{tcolorbox}[colback=gray!5, colframe=gray!50, title=\textbf{Teorema~5.2}]
Bajo las condiciones del Lema~5.2 con fronteras absorbentes en $0$ y en $A$,
si $0 \leq S(0) = i \leq A$, el n\'umero esperado de pasos hasta alcanzar
cualquiera de las dos fronteras es:
\begin{equation}
  \Omega_B(i) = i(A - i).
  \label{eq:eb}
\end{equation}
\end{tcolorbox}

\begin{intuicionmat}
La ecuaci\'on~\eqref{eq:eb} es notable: el tiempo esperado es una par\'abola c\'oncava
en $i$ con m\'aximo en $i = A/2$. Si uno comienza a la mitad del camino,
tarda m\'as en llegar a cualquiera de las dos fronteras. En los extremos
($i = 0$ o $i = A$) ya est\'a en la frontera, as\'i que el tiempo es $0$.
\end{intuicionmat}

\textbf{Ejemplo 5.2.} Para $A = 4$ con $B = \{0, 4\}$:

\begin{center}
\begin{tabular}{c|ccccc}
  \toprule
  $i$ & $0$ & $1$ & $2$ & $3$ & $4$ \\
  \midrule
  $\Omega_B(i)$ & $0$ & $3$ & $4$ & $3$ & $0$ \\
  \bottomrule
\end{tabular}
\end{center}

\subsection{Tiempo de salida condicional (Conditional exit time)}

El \textbf{tiempo de salida condicional (conditional exit time)} $\Omega_A(i)$ es el
promedio de pasos para alcanzar la frontera $A$, condicionado a \emph{no tocar $0$}:

\begin{equation}
  \Omega_A(i) = \frac{\displaystyle\sum_{p_{i\to A}} P_{p_{i\to A}}\,\omega_{p_{i\to A}}}
                     {\displaystyle\sum_{p_{i\to A}} P_{p_{i\to A}}}
             = \frac{\displaystyle\sum_{p_{i\to A}} P_{p_{i\to A}}\,\omega_{p_{i\to A}}}
                    {\mathcal{P}_A(i)}.
  \label{eq:condexitdef}
\end{equation}

Definiendo $z_i = \mathcal{P}_A(i)\,\Omega_A(i)$ y aplicando el Teorema~5.1,
la ecuaci\'on de movimiento resulta en el sistema tridiagonal:

\begin{equation}
  \begin{bmatrix}
    1 & 0 & \cdots & 0 \\
    1 & -2 & \cdots & 0 \\
    \vdots & & \ddots & \vdots \\
    0 & \cdots & -2 & 1 \\
    0 & \cdots & 0 & 1
  \end{bmatrix}
  \begin{bmatrix} z_0 \\ z_1 \\ \vdots \\ z_{A-1} \\ z_A \end{bmatrix}
  =
  \begin{bmatrix} 0 \\ -2/A \\ -4/A \\ \vdots \\ -2(A-1)/A \\ 0 \end{bmatrix}.
  \label{eq:condexitsys}
\end{equation}

\subsection{Teorema 5.3: Tiempo de salida condicional}

\begin{tcolorbox}[colback=gray!5, colframe=gray!50, title=\textbf{Teorema~5.3}]
Bajo las condiciones anteriores, sea $\mathbf{z} = \langle z_0, z_1, \ldots, z_A\rangle$
la soluci\'on de~\eqref{eq:condexitsys}. El tiempo de salida condicional es:
\begin{equation}
  \Omega_A(i) = \frac{i \cdot z_i}{A}, \qquad i = 1, 2, \ldots, A.
  \label{eq:condexit}
\end{equation}
\end{tcolorbox}

\begin{intuicionmat}
La ecuaci\'on~\eqref{eq:condexit} combina la posici\'on inicial $i$, el tama\~no
del objetivo $A$ y el vector auxiliar $z_i$ obtenido del sistema lineal.
Refleja que el tiempo esperado condicionado depende tanto de qu\'e tan cerca
se est\'a del objetivo como de la geometr\'ia global del problema.
\end{intuicionmat}

\textbf{Ejemplo 5.3.} Para $A = 5$ con frontera absorbente en $0$:

\begin{center}
\begin{tabular}{c|ccccc}
  \toprule
  $i$ & $1$ & $2$ & $3$ & $4$ & $5$ \\
  \midrule
  $z_i$ & $8/5$ & $14/5$ & $16/5$ & $12/5$ & $0$ \\
  $\Omega_5(i)$ & $8$ & $7$ & $16/3$ & $3$ & $0$ \\
  \bottomrule
\end{tabular}
\end{center}

%%% =========================================================
\section{Idea Intuitiva de un Proceso Estoc\'astico (Intuitive Idea of a Stochastic Process)}
%%% =========================================================

\subsection{Del modelo determin\'istico al modelo estoc\'astico}

El punto de partida es el \textbf{modelo determin\'istico (deterministic model)} de
crecimiento exponencial. Si la tasa de cambio de una cantidad $P \geq 0$ es
proporcional a $P$:

\begin{equation}
  \frac{dP}{dt} = \mu P,
  \label{eq:det}
\end{equation}

donde $\mu$ es la constante de proporcionalidad (\textbf{tasa de crecimiento (growth rate)}
si $\mu > 0$, \textbf{tasa de decaimiento (decay rate)} si $\mu < 0$).
La soluci\'on es $P(t) = P(0)\,e^{\mu t}$.

\begin{intuicionmat}
La ecuaci\'on~\eqref{eq:det} dice que el crecimiento porcentual por unidad de tiempo
es constante e igual a $\mu$. Es el modelo m\'as simple posible para el
precio de una acci\'on sin aleatoriedad.
\end{intuicionmat}

Dividiendo por $P$:
\begin{equation}
  \frac{dP}{P} = \mu\,dt.
  \label{eq:logdet}
\end{equation}

\begin{intuicionmat}
La ecuaci\'on~\eqref{eq:logdet} expresa el rendimiento logar\'itmico
(\textbf{log-return}) en cada instante: la variaci\'on porcentual del precio
es $\mu\,dt$, determin\'istica.
\end{intuicionmat}

Con el cambio de variable $X = \ln P$:
\begin{equation}
  dX = \mu\,dt.
  \label{eq:logX}
\end{equation}

Su an\'alogo discreto es:
\begin{equation}
  X(t + \Delta t) - X(t) = \Delta X = \mu\,\Delta t.
  \label{eq:disc}
\end{equation}

\begin{intuicion}
El modelo determin\'istico en log-precio es simplemente una l\'inea recta:
$X$ crece a tasa constante $\mu$. No hay sorpresas ni fluctuaciones.
Para capturar la aleatoriedad de los mercados reales hay que a\~nadir un componente
estoc\'astico.
\end{intuicion}

\subsection{Introducci\'on del componente estoc\'astico y volatilidad}

Se introduce una variable aleatoria normal $dz(t)$ con media $0$ y
desviaci\'on est\'andar $1$, independiente en distintos instantes de tiempo.
El \textbf{modelo estoc\'astico (stochastic model)} para el incremento del log-precio es:

\begin{equation}
  \Delta X = \mu\,\Delta t + \sigma\,dz(t)\sqrt{\Delta t},
  \label{eq:stoch}
\end{equation}

donde $\sigma > 0$ se denomina \textbf{volatilidad (volatility)}.

\begin{intuicionmat}
La ecuaci\'on~\eqref{eq:stoch} tiene dos componentes: (1) la parte determin\'istica
$\mu\,\Delta t$, que es la tendencia promedio, y (2) la parte estoc\'astica
$\sigma\,dz(t)\sqrt{\Delta t}$, que es el ``ruido'' aleatorio. El factor
$\sqrt{\Delta t}$ garantiza que la varianza del ruido sea proporcional a $\Delta t$,
lo cual es indispensable para que el proceso tenga propiedades razonables en el
l\'imite continuo.
\end{intuicionmat}

\begin{intuicion}
La volatilidad $\sigma$ mide cu\'anto ``ruido'' tiene el precio. Una acci\'on
con alta volatilidad puede subir o bajar mucho cada d\'ia; una con baja volatilidad
se mueve poco. En mercados financieros, $\sigma$ es uno de los par\'ametros
m\'as importantes para valorar opciones.
\end{intuicion}

\subsection{Media y varianza del incremento acumulado}

Considerando el cambio acumulado sobre el intervalo $[t_i, t_j]$ con
$t_j = t_i + (j-i)\Delta t$:

\begin{equation}
  \Delta X = X(t_j) - X(t_i) = \sum_{k=i}^{j-1}\left(\mu\,\Delta t + \sigma\,dz(t_k)\sqrt{\Delta t}\right).
  \label{eq:accum}
\end{equation}

Como los $dz(t_k)$ son independientes con media $0$:

\begin{align}
  \mathrm{E}[\Delta X] &= \sum_{k=i}^{j-1} \mu\,\Delta t
                        = \mu(j-i)\Delta t = \mu(t_j - t_i).
  \label{eq:mean}
\end{align}

\begin{intuicionmat}
La ecuaci\'on~\eqref{eq:mean} dice que el valor esperado del cambio en log-precio
es la tasa de crecimiento $\mu$ multiplicada por la duraci\'on del intervalo.
El componente estoc\'astico desaparece al tomar la esperanza porque tiene media
cero. El resultado es igual al modelo determin\'istico.
\end{intuicionmat}

La varianza del incremento es:

\begin{align}
  \mathrm{Var}(\Delta X) &= \sigma^2\,\Delta t\,\mathrm{Var}\!\left(\sum_{k=i}^{j-1} dz(t_k)\right)
                          = \sigma^2\,\Delta t\,(j-i)
                          = \sigma^2(t_j - t_i).
  \label{eq:var}
\end{align}

\begin{intuicionmat}
La ecuaci\'on~\eqref{eq:var} revela la raz\'on de incluir $\sqrt{\Delta t}$
en~\eqref{eq:stoch}: hace que la varianza de $\Delta X$ sea \emph{lineal} en
la duraci\'on del intervalo. Sin el factor $\sqrt{\Delta t}$, la varianza
depender\'ia cuadr\'aticamente de $\Delta t$, lo cual es incompatible con las
propiedades del movimiento browniano continuo.
\end{intuicionmat}

\subsection{Proceso de Wiener y paseo aleatorio geom\'etrico}

Se introduce la notaci\'on $dW(t)$ para una variable normal con media $0$ y
varianza $\Delta t$, es decir, $dW(t) = dz(t)\sqrt{\Delta t}$.
El \textbf{proceso de Wiener (Wiener process)} o \textbf{movimiento browniano
(Brownian motion)} es el l\'imite continuo cuando $\Delta t \to 0$.

La ecuaci\'on del paseo aleatorio discreta puede reescribirse como:

\begin{align}
  X(t_{i+1}) &= X(t_i) + \mu\,\Delta t + \sigma\,dW(t_i), \notag \\
  \ln P(t_{i+1}) &= \ln P(t_i) + \mu\,\Delta t + \sigma\,dW(t_i), \notag \\
  P(t_{i+1}) &= P(t_i)\,e^{\mu\,\Delta t + \sigma\,dW(t_i)}.
  \label{eq:geoRW}
\end{align}

\begin{intuicionmat}
La ecuaci\'on~\eqref{eq:geoRW} es el \textbf{paseo aleatorio geom\'etrico
(geometric random walk)}: el precio futuro es el precio actual multiplicado por
un factor aleatorio log-normal. Es el modelo discreto subyacente al famoso
proceso de Black-Scholes en el l\'imite continuo.
\end{intuicionmat}

\begin{intuicion}
El precio del activo se multiplica en cada paso por el factor
$e^{\mu\,\Delta t + \sigma\,dW}$. Cuando $\sigma = 0$ se recupera el crecimiento
determin\'istico $e^{\mu\Delta t}$. Con $\sigma > 0$ aparece la aleatoriedad
log-normal que caracteriza a los modelos de precios de acciones.
\end{intuicion}

%%% =========================================================
\section{Ejemplo del Mercado Accionario (Stock Market Example)}
%%% =========================================================

\subsection{Estimaci\'on de par\'ametros con datos reales}

Se aplica el modelo~\eqref{eq:stoch} a 248 d\'ias de negociaci\'on de las acciones
de \textit{Sony Corporation} (13~agosto~2001 -- 12~agosto~2002).

\textbf{Procedimiento:}
\begin{enumerate}
  \item Se toman los precios de cierre $P_t$ y se calcula $X_t = \ln P_t$.
  \item Se forman los incrementos $\Delta X_t = X_{t+1} - X_t$
        (diferencias logar\'itmicas consecutivas). Con $\Delta t = 1$ (datos diarios):
        \begin{equation}
          \Delta X_t = \mu \cdot 1 + \sigma\,dz(t)\sqrt{1} = \mu + \sigma\,dz(t).
          \label{eq:daily}
        \end{equation}
  \item Si el modelo es correcto, los $\Delta X_t$ deben seguir una distribuci\'on normal.
        La media muestral estima $\mu$ y la desviaci\'on est\'andar muestral estima $\sigma$.
\end{enumerate}

\begin{intuicionmat}
La ecuaci\'on~\eqref{eq:daily} con $\Delta t = 1$ reduce el modelo a una regresi\'on
simple: la distribuci\'on de los retornos diarios es $\mathcal{N}(\mu, \sigma^2)$.
Estimar $\mu$ y $\sigma$ es equivalente a calcular la media y desviaci\'on est\'andar
de la muestra $\{\Delta X_t\}$.
\end{intuicionmat}

\textbf{Resultados emp\'iricos para Sony:}
\begin{center}
\begin{tabular}{ll}
  \toprule
  Par\'ametro & Estimaci\'on \\
  \midrule
  $\mu$ (media diaria) & $\approx -0.000555\ \text{d\'ia}^{-1}$ \\
  $\sigma$ (volatilidad diaria) & $\approx 0.028139\ \text{d\'ia}^{-1}$ \\
  \bottomrule
\end{tabular}
\end{center}

\begin{intuicion}
El valor negativo de $\hat{\mu}$ indica que en promedio el precio de Sony cay\'o
ligeramente cada d\'ia durante ese per\'iodo (coincide con el entorno posterior
al 11-S). La volatilidad $\hat{\sigma} \approx 2.8\%$ diaria es t\'ipica de
una acci\'on de gran capitalizaci\'on. A pesar de la aparente aleatoriedad, los
datos tienen una estructura estad\'istica clara.
\end{intuicion}

%%% =========================================================
\section{Resumen y Conclusiones}
%%% =========================================================

El cap\'itulo establece los fundamentos probabil\'isticos y estoc\'asticos necesarios
para modelar precios de activos financieros. Los resultados principales son:

\begin{itemize}

  \item \textbf{Representaci\'on del paseo aleatorio:}
        $S(N) = S(0) + \sum_{k=1}^N X_k$ con $X_k = \pm 1$ equidistribuidos e i.i.d.
        La distribuci\'on de $S(n)$ es binomial con $p = 1/2$; en $n$ pasos s\'olo son
        alcanzables los estados con la misma paridad que $n$ (Lema~5.1).

  \item \textbf{Principio de reflexi\'on:}
        La f\'ormula de Lema~5.2 permite calcular la probabilidad de alcanzar un umbral $A$
        sin tocar la frontera absorbente en $0$ substrayendo el peso de los caminos
        reflejados.

  \item \textbf{Probabilidades de absorci\'on (Teorema~5.1):}
        Con dos fronteras absorbentes en $0$ y $A$, $\mathcal{P}_A(i) = i/A$ y
        $\mathcal{P}_0(i) = 1 - i/A$. La probabilidad de ruina es proporcional a la
        distancia al objetivo.

  \item \textbf{Tiempo esperado de salida (Teorema~5.2):}
        $\Omega_B(i) = i(A-i)$, obtenido resolviendo la ecuaci\'on de Poisson discreta
        mediante eliminaci\'on gaussiana y sustituci\'on hacia atr\'as.

  \item \textbf{Tiempo de salida condicional (Teorema~5.3):}
        $\Omega_A(i) = i\,z_i/A$, donde $\mathbf{z}$ resuelve un sistema tridiagonal
        con lado derecho que incorpora las probabilidades de absorci\'on.

  \item \textbf{Modelo estoc\'astico continuo:}
        El modelo discreto $\Delta X = \mu\,\Delta t + \sigma\,dz\sqrt{\Delta t}$
        tiene media $\mu(t_j - t_i)$ y varianza $\sigma^2(t_j - t_i)$. La inclusi\'on
        de $\sqrt{\Delta t}$ garantiza la linearidad de la varianza. En el l\'imite
        $\Delta t \to 0$ se obtiene el proceso de Wiener $dX = \mu\,dt + \sigma\,dW$.

  \item \textbf{Aplicaci\'on emp\'irica:}
        Los retornos logar\'itmicos diarios de Sony Corporation presentan una
        distribuci\'on aproximadamente normal con $\hat{\mu} \approx -5.55 \times 10^{-4}$
        d\'ia$^{-1}$ y $\hat{\sigma} \approx 2.81 \times 10^{-2}$ d\'ia$^{-1}$,
        consistente con el modelo estoc\'astico propuesto.

\end{itemize}

\end{document}
